(1) 多方关系为 \(p(v)=p_1(v_1/v)^n\)，代入 \(pv=RT\) 得 \(T(v)=T_1(v_1/v)^{n-1}\)，特别地 \(T_2/T_1=(v_1/v_2)^{n-1}\)。这一步只用了理想气体与给定过程曲线，不要求是绝热过程。
(2) 体积功 \(W=\int_{v_1}^{v_2}K v^{-n}dv=(p_2v_2-p_1v_1)/(1-n)=R(T_2-T_1)/(1-n)\)。内能差 \(\Delta U=C_V(T_2-T_1)\)，第一定律 \(Q=\Delta U+W=[C_V+R/(1-n)](T_2-T_1)\)。所用正号约定为气体对外作功。
(3) 理想气体熵微分为 \(dS=C_VdT/T+Rdv/v\)；积分得到 \(\Delta S=C_V\ln(T_2/T_1)+R\ln(v_2/v_1)\)。用多方关系可写 \(\Delta S=[R-C_V(n-1)]\ln(v_2/v_1)\)，同时可由 \(\int\delta Q_{\rm rev}/T\) 恢复同一式，说明熵是状态函数。
(4) 当 \(n\to1\)，\(T_2\to T_1\)，并可取极限 \(W\to RT_1\ln(v_2/v_1)\)，\(\Delta U\to0\)、\(Q=W\)。理想气体可逆绝热过程满足 \(Q=0\)，故 \(n=1+R/C_V=C_p/C_V\equiv\gamma\)。若实际过程有有限温差换热或摩擦，前述熵公式仍可用于两状态，但不能用 \(\delta Q_{\rm rev}\) 直接代替实际热量。
